% Active article input; prior staged versions remain in Git and validation receipts.
\newcommand{\PaperTitle}{Conditional language generation\\in glyph image space}
\newcommand{\PaperSummary}{Written forms preserved in inscriptions can be absent from a computer font. We investigate language modelling with questions, source passages and replies represented as images of glyphs. In controlled single-paragraph comparisons, a visual writer produced exact two-sentence replies under three new phrasings on 34 of 64 passages. A learned optical correspondence preserved all 192 familiar historical-form replies and raised Kai-script reading from zero to 72 of 192. On Chinese novel passages, matching pictured English questions to their intended answers increased exact stopped replies from 54 to 121 of 748 relative to within-passage exchange. The bilingual writer answered two distinct questions in both languages on three of 374 passages. A separate visual encoder trained on paired questions assigned the correct Chinese and English questions on 364 of 374 novel passages. Conditioning a page reader on these question vectors raised exact English-request excerpts from 17 to 41 of 412 on MLQA while retaining 124 exact Chinese-request excerpts. On two-question novel requests, factorizing an answer into generated clauses with an insertion point for source ink raised correct answers in both orders from zero to 15 of 192 passage--language pairs. A close reading of 48 new requests credited four complete statements for each visual writer and 34 for a text Qwen3-8B comparison. These results show that image-conditioned models can read and write individual grounded answers on a consumer GPU. Reliable open-ended visual language generation and a matched-quality efficiency advantage remain unresolved.}
